\section{Quantitative late phases}\label{sec:late-phases}

Let \(F:\Z_{\ge0}\to[1,\infty)\) be nondecreasing, with
\(F(J)\to\infty\) and
\[
 F(J)^2(1+\log F(J))^2=o(J).
\]
For \(h\ge1\) put \(j(h)=\lfloor\log_2h\rfloor\) and
\(T(h)=\lfloor h F(j(h))\rfloor\). An exceptional set of dyadic
indices will satisfy
\begin{equation}\label{eq:bin-count}
 \#\{j\le J:j\text{ is exceptional}\}
       \ll F(J)^2(1+\log F(J))^2 .
\end{equation}
A conclusion holds on a retained index simultaneously for all integer
\(h\in[2^j,2^{j+1})\). The corresponding retained integers have
logarithmic density one, since each complete dyadic interval has
harmonic weight \(\log 2+O(2^{-j})\).

\subsection{A uniform denominator bound and separation of blocks}

We use the common-weights interval theorem of Stephan
\cite[Theorem~4.2]{Stephan}, in the setting of fixed algebraic real
linear forms and coordinate forms at the primes dividing \(b\).
It extends the interval framework of Evertse--Ferretti
\cite{EF}. For a fixed system of \(n\) forms, suppose the exponents at the real
place are at most \(1\), with one equal to \(1\) and one at most
\(1-\delta/(2n)\). At each finite place they are nonpositive, with
one zero, and the sum of all exponents is at most \(-\delta\).
The inequality for each form is its absolute value bounded by the
sup norm of the integer vector raised to the corresponding exponent,
up to a fixed admissible real factor. Uniformly over the exponent
tuples satisfying these conditions, the theorem confines the
solutions with log height above \(O(\delta^{-1})\) to
\(O(1)\) proper rational subspaces or to log-height intervals
\([z,C\delta^{-1}z)\) of which there are
\(O(\delta^{-2}(1+\log\delta^{-1}))\).
All uses below have real factor \(1\) and real determinant of
absolute value \(1\). This quantitative input is a preprint result.

Suppose a counterexample at \(h\) has log height between
\(c2^{j(h)}\) and \(C(1+F(J))2^{j(h)}\), for \(j(h)\le J\),
with fixed \(c,C>0\). Use \(\delta=c_0/(1+F(J))\) for a
sufficiently small constant \(c_0>0\). A given log-height interval
allows a range of \(2^{j(h)}\) of ratio \(O((1+F(J))^2)\).
It meets \(O(1+\log F(J))\) possible indices \(j(h)\).
The height threshold accounts for \(O(1+\log F(J))\)
additional indices. It follows that the height intervals
contribute \eqref{eq:bin-count}. The points in proper subspaces
will be treated directly.

We will use the following elementary observation for the
three-dimensional cases. Write \(F_J=1+F(J)\). Consider points
\((b^\Delta,y)\) assigned to lengths \(h\) such that
\begin{equation}\label{eq:secant-hyp}
 c h\le\Delta\le C F_J h,\qquad
 c' b^{-C' F_J h}\le |y|\le b^{-c''h}
\end{equation}
for fixed positive constants. A single real affine line contains
such points with at most \(O(1+\log F(J))\) distinct dyadic indices.
Indeed, partition the indices into residue classes modulo
an integer \(O(1+\log F(J))\) so that lengths associated with
increasing indices in a class grow by a factor at least \(K F_J\),
with \(K\) sufficiently large. Except for a bounded number of small
lengths, this makes the horizontal coordinates grow by a factor
at least two and the absolute vertical coordinates fall by a factor
at least four. For three successive such points the horizontal
increment of the second secant is at least that of the first.
The first vertical increment has magnitude at least
\(3|y_1|/4\), while the second has magnitude at most \(5|y_1|/16\).
The secant slopes cannot agree. There are at most two indices
per class, which proves the observation.

\begin{theorem}\label{thm:late-separation}
For every fixed \(\xi>0\), outside a set of indices satisfying
\eqref{eq:bin-count}, all the retained \(h\) satisfy
\begin{equation}\label{eq:uniform-q}
       q_\perp\|q\alpha\|>B^{-\xi}
         \quad(1\le q\le b^{2T(h)}).
\end{equation}
In particular
\(\|k b^m\alpha\|>B^{-\xi}/|k|\) whenever
\(0\le m\le T(h)\) and \(0<|k|\le b^{T(h)}\).

With a possibly enlarged exceptional set satisfying
\eqref{eq:bin-count}, the first \(T(h)\) digits of the original
\(\alpha\) have no two equal \(h\)-digit factors at distinct starts,
with overlap allowed. In this prefix a square has root length less
than \(h\), and two palindromes of different centers have
intersection of length less than \(h\).
\end{theorem}

\begin{proof}
Use the fractional part of \(\alpha\), so \(0<\alpha<1\).
For a counterexample to \eqref{eq:uniform-q} let \(p\) be
the nearest integer to \(q\alpha\), and reduce \(p/q\) to
\(p'/q'\). Its error and its base-external factor have
both decreased, so \(q'_\perp |q'\alpha-p'|\le B^{-\xi}\).
The vector \((q',p')\) is primitive and has norm \(q'\).
Liouville's inequality gives
\(|q'\alpha-p'|\gg (q')^{-(\deg\alpha-1)}\), so the log height
is \(\gg h\), and it is at most \(2T(h)\log b\).
At the real place use \(X_1,\alpha X_1-X_2\) and exponents
\(1,\log|q'\alpha-p'|/\log q'\).
At \(p_0\mid b\) use exponents
\(\log|q'|_{p_0}/\log q',0\). Their total is
\(\log(q'_\perp|q'\alpha-p'|)/\log q'\), which is at most
\(-\delta\) for \(c_0>0\) small in terms of \(\xi\).
Each proper rational line contains at most one primitive
integer vector with positive first coordinate. Its height
can correspond to \(O(1+\log F(J))\) indices by the two height
bounds. The interval estimate then proves \eqref{eq:uniform-q}.
Use \(q=|k|b^m\) for its displayed consequence.

For the repeat assertion fix \(0<\xi<1/4\) in
\eqref{eq:uniform-q}. Let two length-\(h\) factors start after
\(m,n\) places, with \(m<n\) and \(n+h\le T(h)\).
Equality of their labels implies an integer \(p\) with
\[
 |e|=|\alpha(b^n-b^m)-p|<B^{-1}.
\]
If \(\Delta=n-m\le h/2\), the consequence of \eqref{eq:uniform-q}
with \(k=b^\Delta-1\) makes this impossible. Thus \(h/2<\Delta\le
T(h)\). The norm of \((b^n,b^m,p)\) is \(b^n\).
For the real forms
\(X_1,X_2,\alpha(X_1-X_2)-X_3\) use exponents
\(1,m/n,-h/n\). At \(p_0\mid b\) the exponents are
\(\log|b^n|_{p_0}/(n\log b),\,\log|b^m|_{p_0}/(n\log b),\,0\).
The total is \(-h/n\) and the log height
\(n\log b\) is between \(h(\log b)/2\) and \(T(h)\log b\).
Apply the interval estimate. A containing rational plane
\(uX_1+vX_2+wX_3=0\) with \(w=0\) fixes \(\Delta\), leaving
\(O(1+\log F(J))\) possible indices by \(h/2<\Delta\le T(h)\).
When \(w\ne0\), division by \(b^m\) puts
\((b^\Delta,e/b^m)\) on the line
\[
 y=(u/w+\alpha)x+(v/w-\alpha).
\]
Liouville's inequality with denominator \(b^n-b^m\) and the small
error bound give \eqref{eq:secant-hyp}, including constants
independent of the plane. Thus the number of exceptional indices
from the subspaces also satisfies \eqref{eq:bin-count}.

When a square of root length at least \(h\) is wholly in the
prefix, the first \(h\) digits of each copy are such a repeat.
If two palindromes of distinct centers intersect in at least \(h\)
digits, reflect a common \(h\)-digit factor about each center.
Its two reflections are equal factors at distinct starts in the
prefix. This proves the last assertions.
\end{proof}

For example, for any \(\nu>0\) one can take a nondecreasing
modification of
\[
 F(j)=\frac{\sqrt{1+j}}{(1+\log(1+j))^{1+\nu}}.
\]
Then \(T(h)\asymp h\sqrt{\log h}/(\log\log h)^{1+\nu}\), and
the bound \eqref{eq:bin-count} is \(O(J/(1+\log J)^{2\nu})\).
The absence of repeats is a scale-specific quantitative restriction
related to the fixed-positive-proportion exclusion of
Adamczewski--Bugeaud--Luca \cite[Theorem~3.3]{Bilu}.

\subsection{Selectable late phases and growing powers}

\begin{theorem}\label{thm:growing-power}
Fix \(0<\gamma<1/2\). Outside an exceptional set satisfying
\eqref{eq:bin-count}, the following assertions hold.
For fixed \(s\ge2\) and \(0<\theta_0<1\), the counts and seam
construction in Theorem~\ref{thm:selectable-counts} and
Corollary~\ref{cor:phase-palindrome} hold at runs of phases in
\[
 [\,\lfloor\theta_0T(h)\rfloor,\,
       \lfloor\theta_0T(h)\rfloor+(2s+16)h\,].
\]
One may also prescribe any integer \(s=s(h)\) in the range
\(2\le s\le F(j(h))/8\). There is a phase
\begin{equation}\label{eq:growing-phase-window}
 m\in[\,\lfloor T(h)/4\rfloor,\,\lfloor T(h)/4\rfloor+(2s+16)h\,]
\end{equation}
and a positive \(a\equiv1\bmod b\) with
\(a=o(B^{s-\gamma})\) for which \(a\alpha\) has an actual
\(V^s\) of least period \(h\) after \(m\) digits.
The little-oh is uniform for this range of \(s\).
\end{theorem}

\begin{proof}
Choose \(\xi>0\) small, and use the retained indices of
Theorem~\ref{thm:late-separation}. The proof of
Lemma~\ref{lem:phase-residue} uses the Ridout estimate for
\eqref{eq:phase-range} and the bound on \(|U+V\alpha_0|\).
For all the indicated phase windows \(m\le T(h)\). When
\(s\le F(j(h))/8\), the frequencies involved, including \(V\),
satisfy \(|k|,|V|\le b^{T(h)}\) for large \(h\).
The consequence of \eqref{eq:uniform-q} gives both inputs,
with \(\xi<\sigma/4\), and replaces the Ridout constants by uniform
bounds. The parameters \(G,H,\Lambda\) and the determinant
comparisons are independent of \(s\). The final segment
\(j\ge(s+5)h\) still suffices to compare the \(k_j/r_j\).
For fixed \(s\) the Fourier proof above applies directly.
Distinct multipliers of size \(O(A)\) have circle distance
greater than \(1/M\) by the same consequence of \eqref{eq:uniform-q},
using \(\xi<\gamma\).

For growing \(s\), we give the uniform Fourier estimates.
Let \(0<\varepsilon_0<1/100\) be fixed. Convolve a bounded
interval indicator with the \(p\)-fold convolution of the
uniform probability density on
\([-\varepsilon_0/p,\varepsilon_0/p]\).
The resulting smoothings have support in the fixed
\(\varepsilon_0\)-enlargement and total variation at most \(2\).
For derivative order \(k<p\) their \(L^1\) and \(L^\infty\)
derivative bounds are \(O((C p)^k)\). This follows by
differentiating \(k\) distinct factors of the convolution, each
becoming a signed measure of total variation \(p/\varepsilon_0\).
Their Fourier transforms therefore satisfy
\[
 |\widehat f(t)|\ll \min(1,(Cp/|t|)^k).
\]
Finite differences with sampling scale \(A'\) give
\[
 \left|\sum_n f((n+\zeta)/A')e(nt)\right|
   \ll A'\min(1,(Cp/(A'\|t\|))^k)
\]
uniformly in the sampling shift \(\zeta\) if \(k=o(A')\).
In fact the \(k\)-th difference has magnitude
\(O((Cp/A')^k)\) and only \(O(A'+k)\) sampled terms.
The bounded variation gives \(O(1)\) Riemann-sum error.

Fix \(c_1>0\), and choose \(g_p,\psi_p\) of this form
supported in \((0,c_1)\) and \((-1/12,1/12)\), respectively,
with fixed positive integrals. For \(g_p\) choose its smoothing
width small in terms of \(c_1\). Put \(P(x)=1/2-\cos(2\pi x)\).
If \(P(x)\sum_{j}\psi_p(L(Dx+j))>0\), then
\(1/6<x<5/6\) and \(|L\langle Dx\rangle|<1/12\).
Thus \eqref{eq:target-centred} gives a true hit.
Sum over \(a=bn+1\) with weight \(g_p(a/A)\).
At the chosen phases with \(\gamma_0=\gamma+\lambda\), the
nonzero frequencies \(u+vD\) with \(|u|\le1\) and
\(|v|\le LB^{\kappa/2}\) have
\((A/b)\|b(u+vD)\beta\|\gg B^\lambda\), as before.
The zero term is asymptotic to a positive constant times \(A/L\).
The sum \(L^{-1}\sum_v|\widehat\psi_p(v/L)|\) is \(O(p)\)
by the Fourier bound, splitting at a fixed multiple of \(pL\).
The internal error is \(O(Ap(Cp)^k B^{-k\lambda})\).
The tail is \(O(A(Cp)^k B^{-(k-1)\kappa/2})\),
by summing the power decay. Choose
\(k=\lceil C_1(s+1)\rceil\) with \(C_1\) sufficiently large, and
\(p=k+2\). Since \(\log p=o(h)\), both errors are \(o(A/L)\)
uniformly over the stated range of \(s\).

All weights are bounded, so there are \(\gg_{c_1} A/L\)
positive contributions. The hits use distinct root cylinders by
the circle distance bound just proved. Since \(A/L\asymp
B^{1-\gamma}\), one can discard every root with a period
at most \(h/2\), leaving powers of least period \(h\)
by the cyclic-shift argument of Corollary~\ref{cor:phase-palindrome}.
For every fixed \(c_1>0\) there is such a hit with \(0<a<c_1A\).
Taking the minimum proves the uniform little-oh statement.
\end{proof}

For \(s=\lfloor F(j(h))/8\rfloor\) the number of copies is unbounded
and \(sh\) is comparable to \(T(h)\). In \eqref{eq:growing-phase-window}
the complete block lies in the first \(T(h)\) digits for sufficiently
large \(h\). Theorem~\ref{thm:late-separation} excludes the same
root length in the original digits within that window on the retained
indices. The two results identify where a transfer from changing
multipliers needs additional information.

